%% Pourquoi la mission existe : le probleme que l'entreprise a voulu
%% traiter, ferme par une question a laquelle le rapport repond.

\subsection{Why ALTEN Opened This Subject}
\label{sec:pourquoi}

Every component the programme had built by the start of this internship answers a question of
the form \emph{what is the state of this part of the chain?} None answers the question that
comes next. Once the state is known, a coordinator still has to decide where to spend a finite
amount of attention this week, and that decision runs on a signal none of those modules
produces: how urgent each node is. In practice it comes from people. A buyer says a supplier is
critical; a plant says a line is at risk. The twin records the consequence of those judgments
without ever evaluating the judgments themselves.

That is a gap with a cost attached, because the signal is not neutral. Zhu, Yang and Hsee
demonstrate across five controlled experiments that people systematically favour imminent
low-value tasks over distant high-value ones, and that the bias persists even when they are
told about it~\autocite{zhu2018mere}; field studies find the same pattern under
workload~\autocite{rusou2020psychology,diwa2020task}. A twin that consumes a human urgency
declaration as ground truth inherits that bias and propagates it across the network.

\subsection{The Cost of the Two Errors}

The two directions of error are not symmetric, which is what makes the problem worth
instrumenting rather than tolerating. Figure~\ref{fig:asymetrie} sets them side by side.

\begin{figure}[H]
\centering
\begin{tikzpicture}[
    err/.style = {bloc, text width=5.4cm, minimum height=0.82cm, font=\scriptsize},
    errH/.style = {err, fill=altenOchrePale!55!white, draw=altenOchreDark, line width=1pt},
    node distance = 3mm]

\node[blocFort, text width=5.2cm, font=\small] (dec)
     {\textbf{the urgency an operator declares}};

\node[err, below left=11mm and 3mm of dec.south, anchor=north east, fill=altenBlueBg2]
     (f1) {\textbf{too high: false urgency}\\[1pt] $U_d > U_r$};
\node[errH, below right=11mm and 3mm of dec.south, anchor=north west]
     (h1) {\textbf{too low: hidden risk}\\[1pt] $U_d < U_r$};

\node[err, below=of f1] (f2) {consumes expedited transport, premium capacity and the
                              coordinator's attention};
\node[err, below=of h1] (h2) {raises no alert, enters no priority list, consumes no
                              attention};

\node[err, below=of f2] (f3) {the bill lands on another flow, often on another actor
                              entirely};
\node[err, below=of h2] (h3) {surfaces only as a missed delivery, once the window to act on
                              it has closed};

\node[blocPale, text width=5.4cm, minimum height=0.82cm, font=\scriptsize, below=of f3]
     (f4) {\textbf{visible:} somebody notices the capacity is gone};
\node[errH, below=of h3] (h4) {\textbf{silent:} nothing in the system records that it
                               happened};

\draw[fleche] (dec.south) -- ++(0,-4.5mm) -| (f1.north);
\draw[fleche] (dec.south) -- ++(0,-4.5mm) -| (h1.north);

\coordinate (mid) at ($(f4.south)!0.5!(h4.south)$);
\node[blocAlerte, below=7mm of mid, text width=11.1cm, font=\scriptsize] (pen)
     {Only the second error is silent, so the adequacy score prices it at roughly twice the
      first. It is the rule clinical triage has applied for decades.};
\draw[fleche] (f4.south) -- (f4.south |- pen.north);
\draw[fleche] (h4.south) -- (h4.south |- pen.north);

\end{tikzpicture}
\caption[The two declaration errors and why only one of them is silent]{The two directions in
which a declaration can be wrong.}
\label{fig:asymetrie}
\end{figure}

Both costs are documented. In a shared or mutualised logistics setting the price of
over-declaring lands on another actor
entirely~\autocite{montreuil2011toward,li2020exploring}, and clinical triage, where the
asymmetry has been studied for decades, reaches the same verdict on the other direction:
under-triage costs a life where over-triage costs a
queue~\autocite{kim2020admission,mackway2013emergency}. That asymmetry is the argument behind
the penalty of Section~\ref{sec:realisation_modele}.

\subsection{The Problem, Stated}
\label{sec:question}

\begin{definitionBox}
\noindent\textbf{Problem statement.}

\medskip

A supply chain digital twin can compute the physical state of every node, and it can record
the urgency a human declares for each of them. It does not compare the two. The declared
signal is known to be biased, the computed signal is available, and the difference between
them is exactly the quantity that would tell a coordinator which node is being mismanaged
rather than merely which node is struggling.

\medskip

\noindent\emph{Can the gap between declared and computed urgency be measured reliably
enough, inside a working digital twin, to be acted upon rather than merely observed?}
\end{definitionBox}

The word \emph{reliably} carries the weight. Producing a number that expresses the difference
between two other numbers is trivial. Establishing that the number means something, and that it
predicts outcomes better than the raw indicators a coordinator already has, is the actual
engineering research problem.
